\scanpage{53}

Supposons $k\geq0$, et le lemme \unsure{prouvé} pour $k'<k$. On sait déjà que les $\unsure{p_i|\rho_i}:H^*(X)\mathrel{-}H^i(X)$ sont $\in\unsure{\operatorname{Alg}|A}$ pour $i<k$, donc (transp.) pour $i>2n-k$. Soit
\[
T_k=\operatorname{id}-\sum_{i<k}\unsure{p_i|\rho_i}-\sum_{j>2n-k}\unsure{p_j|\rho_j},
\]
donc $T_k\in A$, plus gén. $T_{k'}\in A$ si $k'\unsure{\leq|<}k$. Alors $\xi^{n-k}T_k$ est nul sur $H^{\unsure{j|i}}$, $j\neq k$, et induit
\diagram{\begin{tikzcd}[ampersand replacement=\&]
H^k\arrow[r,"\xi^{n-k}","\sim"']\&H^{2n-k}
\end{tikzcd}}
, donc $u_{\unsure{k|n}}\xi^{n-k}T_k=\unsure{p_k|\rho_k}$.

D'autre part, si $i+2j=k$, \inserted{$j>0$,} la projection sur $\xi^jP^i$ est aussi $\in A$, comme on voit \deleted{\unsure{ainsi}} par récurrence ascendante sur $i$ (donc descendante sur $j$) ainsi : si $i<0$ rien à prouver. Si c'est fait pour les $i'<i$, on forme $\xi^{n-k+j}T_k$, \deleted{$\unsure{\xi}$} qui induit zéro sur $H^\alpha$ pour $\alpha\neq k$, et induit
\diagram{\begin{tikzcd}[ampersand replacement=\&]
H^k\arrow[r,"\xi^{\unsure{n-k-j|n-k+j}}"]\&H^{2n-k+2j}
\end{tikzcd}}
, \unsure{qui est nul} \deleted{\unsure{sur}} sur les facteurs $\xi^{j'}P^{i'}$ pour $j'<j$, \inserted{de noyau $N$,} et sur les facteurs $\xi^{j'}P^{i'}$ avec $j'\geq j$ (i.e. $i'\leq i$), \deleted{\unsure{coïncide à} $\unsure{1}$} dont la somme n'est autre que $\xi^jH^{k-2j}$, son composé avec l'isom
\diagram{\begin{tikzcd}[ampersand replacement=\&]
H^{k-2j}\arrow[r,"\xi^j","\sim"']\&\xi^jH^{k-2j}
\end{tikzcd}}
est $\xi^{n-(k-2j)}$. Donc l'hom $u_{k-2j}\xi^{n-k+j}T_k$ \unsure{est}, nul sur $N$, et sur $\xi^jH^{k-2j}$ l'inverse de
\diagram{\begin{tikzcd}[ampersand replacement=\&]
H^{k-2j}\arrow[r,"\sim"]\&\xi^jH^{k-2j}
\end{tikzcd}}
. Donc

\mspage{52}
